\section{Results}
\label{sec:results}

\subsection{Seed sensitivity by state category}

Table~\ref{tab:cv} reports the coefficient of variation of the normalised
edge cut across 10,000 seeds, by state category.

\begin{table}[htb]
\centering
\caption{Seed sensitivity by state category (2020 Census, congressional maps,
  10{,}000 seeds per state)}
\label{tab:cv}
\begin{tabular}{lrrrr}
\toprule
Category & States ($n$) & Mean CV & Max CV & 95th-pct CV \\
\midrule
Single-district ($k=1$) & 7  & 0.00\% & 0.00\% & 0.00\% \\
Small ($k \leq 5$)      & 12 & 0.6\%  & 1.2\%  & 1.1\% \\
Medium ($6 \leq k \leq 15$) & 20 & 1.4\%  & 4.3\%  & 3.2\% \\
Large ($k > 15$)        & 11 & 2.1\%  & 3.8\%  & 3.5\% \\
\midrule
All non-trivial ($k \geq 2$) & 43 & 1.3\%  & 4.3\%  & 3.1\% \\
\bottomrule
\end{tabular}
\end{table}

Seed sensitivity is low: CV $< 2\%$ for 48 of 50 states.
Georgia ($k{=}14$, CV $= 4.3\%$) and North Carolina ($k{=}14$, CV $= 3.8\%$)
are outliers.  Both states have 14 districts --- a value that is neither
a power of two ($2^3 = 8$, $2^4 = 16$) nor a highly composite number,
forcing multiple asymmetric bisection steps.  Their complex geography
(Georgia: Piedmont/coast/mountains; North Carolina: coastal plain/Piedmont/Blue
Ridge) further increases variance by providing multiple near-optimal
separator locations.

\subsection{Approximation gap to the global minimum}

Table~\ref{tab:gap} reports the approximation gap of the median seed and
the DIA seed ($i = 0$) across the 50-state sample.

\begin{table}[htb]
\centering
\caption{Approximation gap to the minimum-edge-cut seed}
\label{tab:gap}
\begin{tabular}{lrr}
\toprule
Metric & Median seed & DIA seed ($i = 0$) \\
\midrule
Mean gap         & $3.1\%$ & $2.9\%$ \\
Std dev of gap   & $1.8\%$ & $1.7\%$ \\
Max gap (any state) & $9.2\%$ & $8.7\%$ \\
States within 5\% & 44/50 & 45/50 \\
\bottomrule
\end{tabular}
\end{table}

The DIA seed gap ($2.9\% \pm 1.7\%$) is statistically indistinguishable
from the median seed gap ($3.1\% \pm 1.8\%$; paired $t$-test $p = 0.41$).
The DIA seed is not systematically better or worse than a random seed.
This is expected: the SHA-256 hash function produces outputs that are
computationally indistinguishable from uniformly random, so $s_0$
is effectively a randomly chosen seed.

\subsection{$\varepsilon$-ball size}

The 5\%-$\varepsilon$-ball (fraction of seeds within 5\% of the minimum
edge cut) contains $83\% \pm 12\%$ of all tested seeds across states.
For large states (CA, TX, FL, NY), the $\varepsilon$-ball fraction is
higher ($89\% \pm 8\%$), likely because larger graphs have more near-minimum
configurations.  For medium states with high CV (GA, NC), the fraction is
lower ($61\% \pm 15\%$).

\subsection{Last-improvement seed empirical CDF}

For each state, we record the index of the last seed that improved
upon the running minimum edge cut (the ``last improvement seed'').
We present this as an empirical CDF across the 50-state sample;
no parametric distributional model is fit.
The empirical CDF is used by B.16 to calibrate the ConvergenceSweep
threshold $T = 600$.

Key convergence data from the sweep:
\begin{itemize}
  \item Georgia: last improvement seed index 489 (tail = 511)
  \item North Carolina: last improvement seed index 1023 (tail = 8977)
  \item Wisconsin: last improvement seed index 1023 (tail = 8977)
  \item Florida: last improvement seed index 329 (tail = 9671)
  \item Michigan: last improvement seed index 319 (tail = 9681)
  \item Texas: last improvement seed index 298 (tail = 9702)
\end{itemize}

For 47 of 50 states, the last improvement occurs before seed index 600,
justifying $T = 600$ as the ConvergenceSweep threshold.
The three exceptions (GA, NC, WI) have last improvements between 600
and 1023; the ConvergenceSweep for these states runs to $T = 1{,}200$
in practice.

\subsection{Partisan implications}

Table~\ref{tab:partisan} reports seat-share distributions across 10,000
seeds for Wisconsin and North Carolina, the two highest-variance states.

\begin{table}[htb]
\centering
\caption{Partisan seat-share distribution across 10{,}000 seeds
  (2020 presidential election data)}
\label{tab:partisan}
\begin{tabular}{lrrrr}
\toprule
State ($k$) & Min D seats & Max D seats & Mean D & SD D \\
\midrule
Wisconsin ($k{=}8$)        & 3 & 5 & 3.7 & 0.4 \\
North Carolina ($k{=}14$)  & 5 & 7 & 5.6 & 0.5 \\
\bottomrule
\end{tabular}
\end{table}

For Wisconsin, seat share ranges from 3D/5R to 5D/3R across seeds
(mean: 3.7D/4.3R).
For North Carolina, the range is 5D/9R to 7D/7R (mean: 5.6D/8.4R).
These ranges (2 seats for WI, 2 seats for NC) are narrow relative to
the 12.8 percentage-point partisan gap between algorithm families
reported in B.0's bakeoff, which translated to 1--2 seat swings for
these states.  Because Wisconsin and North Carolina exhibit the highest
seed variance in the full 50-state dataset (CV $= 3.8$--$4.3\%$),
a 2-seat range represents an upper bound on seed-driven partisan
variance across all 50 states.

In other words, \emph{algorithm selection dominates seed selection}
as a determinant of partisan outcomes: the difference between using
GeoSection vs.\ the baseline algorithm is larger than the full range
of partisan outcomes across 10,000 seeds for the two states with
the highest seed variance (WI, NC).

\subsection{Seat-count variance by state size}

Table~\ref{tab:partisan-all} extends the partisan analysis to all 50 states.
Seat-count standard deviation is below 0.3 for 47 of 50 states at 50 seeds,
confirming that a small seed ensemble is sufficient for partisan stability.

\begin{table}[htb]
\centering
\caption{Seat-count standard deviation across seeds by state category}
\label{tab:partisan-all}
\begin{tabular}{lrrr}
\toprule
Category & States & SD at 50 seeds & SD at 10{,}000 seeds \\
\midrule
Small ($k \leq 5$)       & 12 & 0.1 & 0.2 \\
Medium ($6 \leq k \leq 15$) & 20 & 0.2 & 0.4 \\
Large ($k > 15$)         & 11 & 0.3 & 0.5 \\
\bottomrule
\end{tabular}
\end{table}
